\chapter*{INTRODUCTION}
\phantomsection
\addcontentsline{toc}{chapter}{INTRODUCTION}

Cervical cancer remains one of the most preventable causes of cancer morbidity and mortality among women. Persistent infection with high-risk human papillomavirus (HPV) types is the necessary cause of nearly all cervical cancer, and prophylactic HPV vaccination can prevent the majority of HPV-related cervical precancer and cancer when delivered before or around the period of sexual debut \cite{schiffman2016,who2022hpv}. The global strategy for cervical cancer elimination therefore places HPV vaccination alongside screening and treatment as a central public health intervention \cite{who2020strategy}.

Despite this strong biological and programmatic rationale, HPV vaccination is not only a biomedical issue. In countries where vaccination depends partly on private payment, geographic availability, and individual information-seeking, uptake can become strongly patterned by household wealth, education, residence, ethnicity, and media access. These inequalities matter because they can reproduce the same social gradient already seen in cervical cancer screening and treatment access: women with fewer resources may be less likely to receive both primary prevention and early detection.

Viet Nam provides an important setting for examining this issue. Previous Vietnamese evidence has largely come from local, clinic-based, student-based, or knowledge-and-acceptability studies; nationally representative evidence quantifying wealth-related inequality in awareness, self-reported any-dose uptake, and the awareness--vaccination gap using concentration-index methods remains limited.  The Viet Nam Multiple Indicator Cluster Survey 2020--2021 (MICS6) collected nationally representative information on women, households, media exposure, reproductive health, and HPV vaccination awareness and uptake \cite{unicef2022}. The survey predates the planned inclusion and public financing of cervical cancer prevention vaccine in the Expanded Programme on Immunization from 2026 under Resolution 104/NQ-CP and the 2026--2028 EPI implementation plan \cite{govvn2022nq104,baochinhphu2025epi}. Among young women aged 15--29 years, the dataset therefore permits analysis of both early adult vaccination experience and a pre-program equity baseline before public immunization financing is expanded.

This thesis focuses on three policy-relevant outcomes: HPV vaccine awareness, self-reported any-dose uptake, and remaining unvaccinated despite awareness. The last outcome is especially important because it distinguishes awareness from uptake and can signal barriers beyond information alone. A woman may know about the vaccine but still remain unvaccinated because of possible cost, distance, limited service availability, family or social constraints, low perceived priority, or other unmeasured factors.

The specific objectives were:
\begin{enumerate}
    \item To estimate the survey-weighted prevalence of HPV vaccine awareness, self-reported any-dose HPV vaccine uptake, and the awareness--vaccination gap among Vietnamese women aged 15--29 years.
    \item To describe how these outcomes vary by sociodemographic and access-related characteristics, using survey-corrected bivariate tests.
    \item To estimate adjusted associations between explanatory variables and each outcome using survey-weighted Poisson regression.
    \item To quantify socioeconomic inequality using standard and Erreygers-corrected concentration indices.
    \item To decompose the Erreygers concentration index to describe measured components that account for observed inequality.
\end{enumerate}

\clearpage
